\documentclass[10pt]{article}
\usepackage[T1]{fontenc}
\usepackage[utf8]{inputenc}
\usepackage[french]{babel}
\usepackage{newpxtext,newpxmath}
\usepackage[a4paper,top=1.5cm,bottom=1.35cm,left=1.6cm,right=1.6cm]{geometry}
\usepackage{microtype}
\usepackage{xcolor}
\usepackage{titlesec}
\usepackage[most]{tcolorbox}
\usepackage{enumitem}
\usepackage{array}
\usepackage{booktabs}
\usepackage{fancyhdr}

\definecolor{navy}{HTML}{1F3A5F}
\definecolor{grey}{HTML}{54606E}
\definecolor{green}{HTML}{2C6E49}
\definecolor{amber}{HTML}{9C6A16}
\definecolor{red}{HTML}{A6322E}
\definecolor{plum}{HTML}{7A4EA8}
\definecolor{boxgray}{HTML}{F3F5F8}
\definecolor{boxred}{HTML}{FCF2F1}
\definecolor{boxgreen}{HTML}{F1F7F3}
\definecolor{boxplum}{HTML}{F7F3FB}
\definecolor{boxcode}{HTML}{F7F8FA}
\definecolor{boxedge}{HTML}{D9E0EA}

\setlength{\parindent}{0pt}
\setlength{\parskip}{2.6pt}
\linespread{1.0}

\pagestyle{fancy}\fancyhf{}
\lhead{\small\color{grey}Scella · architecture}
\rhead{\small\color{grey}M0 --- Le profil servi --- 20 août 2026}
\cfoot{\small\color{grey}\thepage}
\renewcommand{\headrulewidth}{0.3pt}

\titleformat{\section}{\color{navy}\large\bfseries}{}{0pt}{}
\titlespacing{\section}{0pt}{11pt}{3pt}
\titleformat{\subsection}{\color{navy}\bfseries}{}{0pt}{}
\titlespacing{\subsection}{0pt}{7pt}{1.5pt}

\newtcolorbox{these}{colback=boxgray,colframe=boxedge,boxrule=0.5pt,arc=2.5pt,
  left=8pt,right=8pt,top=4.5pt,bottom=4.5pt}
\newtcolorbox{alerte}{colback=boxred,colframe=red!55,boxrule=0.8pt,arc=2.5pt,
  left=8pt,right=8pt,top=4.5pt,bottom=4.5pt}
\newtcolorbox{acquis}{colback=boxgreen,colframe=green!45,boxrule=0.8pt,arc=2.5pt,
  left=8pt,right=8pt,top=4.5pt,bottom=4.5pt}
\newtcolorbox{piste}{colback=boxplum,colframe=plum!45,boxrule=0.8pt,arc=2.5pt,
  left=8pt,right=8pt,top=4.5pt,bottom=4.5pt}
\newtcolorbox{extrait}{colback=boxcode,colframe=boxedge,boxrule=0.5pt,arc=1.5pt,
  left=7pt,right=7pt,top=3.5pt,bottom=3.5pt}

\newcommand{\best}[1]{\textcolor{green}{\textbf{#1}}}
\newcommand{\pire}[1]{\textcolor{red}{\textbf{#1}}}
\newcommand{\sz}[1]{\textcolor{grey}{\small\textsf{#1}}}
\newcommand{\code}[1]{\texttt{\small #1}}
\newcommand{\fl}{$\rightarrow$}

\begin{document}

\begin{center}
{\Huge\bfseries\color{navy} Scella}\\[4pt]
{\large\color{grey}\itshape M0 --- Le profil servi · architecture}\\[4pt]
{\small\color{grey} 20 août 2026 \; · \; premier module du chemin critique \; · \; trois jours}
\end{center}
\vspace{2pt}

\begin{these}
\textbf{Le principe tient en une phrase.} \emph{Le binaire dit ce qu'il sert ; le profil déclare ce
qu'on attend ; un vérificateur compare les deux --- et aucun des trois ne peut diverger sans casser un
test.} Un profil n'est pas un fichier \code{.env} documenté : ce serait un second dépôt de vérité, et
il dériverait au premier interrupteur ajouté. C'est un \textbf{triplet contraint}.

\textbf{Et le manque n'est pas la rigueur.} La passerelle refuse déjà de démarrer sur un master de
clés invalide (\code{exit 78}), garde la frontière derrière quatre préconditions, et se refuse à
imprimer un « OK » lisse quand la capacité PII manque. Ce qui manque est la \textbf{complétude} et la
\textbf{lisibilité machine} : les \code{println!("→ …")} sont une auto-description en prose, partielle,
que rien ne contraint.
\end{these}

\section{1 · Le problème exact --- et une correction}

\begin{alerte}
\textbf{La feuille de route disait « sept interrupteurs ». C'est faux}, et l'écart est instructif.
Mesuré : \textbf{25} variables d'environnement lues par \code{efe-swarm-gateway}, \textbf{7} par
\code{audit-console}. Trente-deux, sur deux binaires.

Deux corrections de détail au passage : \code{EFE\_K} \textbf{n'existe pas} --- mon relevé attrapait
\code{EFE\_Kovant\_P1\_*} dans des commentaires ; et \code{EFE\_META\_PROVENANCE\_TENANT} appartient à
\textbf{C2}, pas au périmètre C1/C3 de ce module.
\end{alerte}

Trente-deux variables, et \textbf{aucun endroit qui dise ce que la configuration résultante sert
réellement}. C'est cela, et rien d'autre, que M0 corrige.

\section{2 · Ce qui existe déjà --- et qu'il ne faut pas réinventer}

\renewcommand{\arraystretch}{1.22}
\begin{center}\small
\begin{tabular}{@{}>{\raggedright\arraybackslash}p{5.4cm} >{\raggedright\arraybackslash}p{10.4cm}@{}}
\toprule
\textbf{Acquis} & \textbf{Ce qu'il fait déjà} \\
\midrule
\code{exit 78} \sz{(EX\_CONFIG)} & master de clés fourni mais invalide \fl{} \textbf{refus de démarrer}, non-retryable \\
\code{frontier\_boot\_guard(…)} & fonction \textbf{pure}, testée, \textbf{quatre préconditions} fail-closed sur l'egress payant \\
Le préflight honnête & \textbf{refuse d'imprimer un « OK » lisse} quand la capacité PII est absente --- il répète le risque \\
\code{/readyz} & readiness \emph{réelle} : 503 si la frontière est annoncée mais non câblée \\
\bottomrule
\end{tabular}
\end{center}

\begin{acquis}
\textbf{M0 généralise cette discipline au reste de la surface. Il ne l'invente pas.} C'est important
pour le périmètre : le module est un travail d'\emph{extension} et de \emph{mise en ordre}, pas de
conception d'un système de configuration.
\end{acquis}

\section{3 · Les quatre pièces}

\subsection{3.1 · Le registre --- source unique de vérité}

Chaque interrupteur déclaré \textbf{une seule fois}, avec ce qu'il achète \emph{et ce que son absence
coûte} :

\begin{extrait}
\code{Switch \{} \\
\code{\ \ \ \ name:\ \ \ \ \ \ \ \ \ \ "EFE\_NANO\_AUDIT\_HMAC\_MASTER",} \\
\code{\ \ \ \ buys:\ \ \ \ \ \ \ \ \ \ "provenance non-forgeable",} \\
\code{\ \ \ \ absence\_costs:\ "repli sur la clé de TEST PUBLIQUE --- le reçu devient forgeable",} \\
\code{\ \ \ \ class:\ \ \ \ \ \ \ \ \ Class::Degraded,} \\
\code{\ \ \ \ surface:\ \ \ \ \ \ \ Surface::Receipt,\ \ // où l'absence se voit, côté servi} \\
\code{\}}
\end{extrait}

\begin{alerte}
\textbf{Contrainte dure : le registre est le seul endroit du crate qui lit l'environnement.} Sans
cela il devient une documentation parallèle --- c'est-à-dire exactement ce que M0 existe pour
supprimer.
\end{alerte}

\subsection{3.2 · La taxonomie --- trois classes, une règle d'or}

\begin{center}\small
\begin{tabular}{@{}>{\raggedright\arraybackslash}p{3.4cm} >{\raggedright\arraybackslash}p{4.6cm} >{\raggedright\arraybackslash}p{7.6cm}@{}}
\toprule
\textbf{Classe} & \textbf{Absence \fl{}} & \textbf{Réservée à} \\
\midrule
\textbf{A · Fatale} & \pire{\code{exit 78}}, refus de démarrer & ce dont l'absence rendrait une affirmation \textbf{fausse} \sz{(la frontière l'utilise déjà)} \\
\textbf{B · Dégradée annoncée} & démarre, capacité \textbf{éteinte}, \textbf{dite en clair} & tout le reste qui touche à la sûreté \\
\textbf{C · Réglage} & défaut sûr, silence acceptable & capacité, adresses, tailles \\
\bottomrule
\end{tabular}
\end{center}

\begin{acquis}
\textbf{La règle d'or : aucune classe B silencieuse.} Une capacité éteinte se dit à \textbf{trois}
endroits --- le rapport au démarrage, la surface servie, et le reçu. Le précédent existe déjà dans le
dépôt : \emph{racine vide $\Leftrightarrow$ aucun scellement}. M0 ne fait que le généraliser.
\end{acquis}

\subsection{3.3 · Le rapport --- dérivé du code qui décide}

Deux formes, \textbf{une seule source} : la forme \textbf{humaine} au démarrage --- les lignes
\code{→} d'aujourd'hui, mais \emph{complètes} --- et la forme \textbf{machine}, via
\code{-{}-profile-report}, qui imprime le JSON et \textbf{sort} sans écouter.

Par capacité, le rapport dit : \best{SERVIE} · \textcolor{amber}{\textbf{DÉGRADÉE}} ·
\pire{REFUSÉE} --- et pour dégradée, \textbf{ce que ça coûte}.

\begin{piste}
\textbf{Pourquoi un drapeau et non une route \code{/profile}.} Le rapport dit ce qui est \emph{éteint} :
« capacité PII absente » est un renseignement pour un attaquant. Un drapeau de ligne de commande
n'ajoute \textbf{aucune surface d'attaque} ; une route en ajoute une qu'il faudrait ensuite protéger,
authentifier et tester. Pour trois jours, le drapeau gagne --- et il suffit au vérificateur.
\end{piste}

\subsection{3.4 · Les profils nommés}

\begin{center}\small
\begin{tabular}{@{}>{\raggedright\arraybackslash}p{2.5cm} >{\raggedright\arraybackslash}p{6.4cm} >{\raggedright\arraybackslash}p{6.6cm}@{}}
\toprule
\textbf{Profil} & \textbf{Ce qu'il sert} & \textbf{Ce qu'il refuse} \\
\midrule
\code{demo} & tout local, clé de dev, aucun egress & prétendre que le reçu vaut quelque chose --- le rapport dit \pire{« forgeable »} \\
\code{souverain} & scellement réel, ancre, bundle, SLM local & tout appel frontière \\
\code{frontiere} & souverain \textbf{+} frontière & démarrer si l'une des 4 préconditions manque \\
\bottomrule
\end{tabular}
\end{center}

Chaque profil est livré avec son \code{profiles/<nom>.expected.json} --- le rapport que le binaire
\textbf{doit} produire.

\section{4 · Le garde-fou --- ce qui empêche la dérive}

C'est la pièce qui sépare M0 d'un document qui pourrit en trois semaines. Deux tests, tous deux dans
l'idiome du dépôt : \emph{les invariants sont des tests qui cassent le build}.

\begin{center}\small
\begin{tabular}{@{}>{\raggedright\arraybackslash}p{4.6cm} >{\raggedright\arraybackslash}p{11.2cm}@{}}
\toprule
\textbf{Test} & \textbf{Ce qu'il interdit} \\
\midrule
\textbf{T1} \code{aucun\_env\_var\_hors\_registre} & toute occurrence de \code{std::env::var(} \textbf{hors} du module \code{profile} fait échouer le test. C'est ce qui rend le registre réellement unique \\
\textbf{T2} \code{le\_rapport\_couvre\_le\_registre} & chaque interrupteur apparaît dans le rapport, et chaque profil produit exactement son \code{expected.json}. \textbf{Ajouter un interrupteur sans le documenter casse les tests} \\
\bottomrule
\end{tabular}
\end{center}

\section{5 · Le critère d'acceptation, exécutable}

\begin{extrait}
\code{git clone <dépôt> \&\& cd efe} \\
\code{set -a \&\& . ./profiles/souverain.env \&\& set +a} \\[3pt]
\code{cargo run -q -p efe-swarm-gateway -- -{}-profile-report | jq -e '.provenance == "SERVIE"'} \\
\code{cargo run -q -p efe-swarm-gateway \&} \\
\code{curl -s localhost:8080/v1/swarm/run -d @examples/run.json \textbackslash} \\
\code{\ \ | jq -e '.provenance.root | length == 64'}
\end{extrait}

\textbf{Fini quand les deux \code{jq -e} sortent 0.} Le premier prouve que le binaire \emph{déclare}
servir ; le second qu'il \emph{sert}. Ensemble, ils ferment l'écart entre le déclaré et le servi ---
ce qui est l'objet même du module.

\section{6 · Le périmètre --- ce que M0 ne fait pas}

Trois jours, donc l'exclusion est structurante.

\begin{itemize}[leftmargin=1.2em,itemsep=2pt,topsep=2pt]
\item \textbf{Pas de langage de configuration} --- des fichiers \code{.env} et un registre suffisent.
\item \textbf{Pas de gestionnaire de secrets} --- le master reste une variable d'environnement ; sa
      provenance est un problème d'exploitation, pas d'architecture.
\item \textbf{Pas de rechargement à chaud} --- un profil se change par redémarrage.
\item \textbf{Pas de configuration par tenant} --- l'isolation existe déjà, c'est un autre sujet.
\item \textbf{Pas de crate partagé} --- les deux binaires ont des surfaces \textbf{disjointes} (25 et 7,
      \emph{zéro recouvrement}). Un module \code{src/profile.rs} par binaire, format de rapport commun
      par convention. Créer une infrastructure partagée pour deux consommateurs serait de
      l'anticipation.
\end{itemize}

\section{7 · Les risques}

\renewcommand{\arraystretch}{1.22}
\begin{center}\small
\begin{tabular}{@{}>{\raggedright\arraybackslash}p{6.2cm} >{\raggedright\arraybackslash}p{9.6cm}@{}}
\toprule
\textbf{Risque} & \textbf{Mitigation} \\
\midrule
Le registre devient une \emph{seconde} vérité & \textbf{T1} --- aucun \code{env::var} hors registre \\
Le rapport \textbf{affirme plus} que le code ne fait & le rapport est calculé depuis les mêmes valeurs que le constructeur d'état consomme, jamais depuis une liste parallèle \\
Les profils pourrissent & un \code{expected.json} par profil, comparé en intégration continue \\
Le rapport renseigne un attaquant & drapeau de ligne de commande, jamais de route \\
\bottomrule
\end{tabular}
\end{center}

\section{8 · Le découpage --- trois jours}

\begin{center}\small
\begin{tabular}{@{}c >{\raggedright\arraybackslash}p{14.4cm}@{}}
\toprule
& \textbf{Ce qui se fait} \\
\midrule
\textbf{J1} & Le registre : inventorier les 25 $+$ 7, classer A/B/C, router \emph{toutes} les lectures. \textbf{T1 écrit en premier} --- il révèle immédiatement ce qui a été oublié. \\
\textbf{J2} & Le rapport, humain et JSON, dérivé du registre. \code{main.rs} consomme le rapport au lieu des \code{println!} épars. \textbf{T2} et les trois \code{expected.json}. \\
\textbf{J3} & Les trois profils, le critère d'acceptation de bout en bout, et le même travail --- plus court --- sur \code{audit-console}. \\
\bottomrule
\end{tabular}
\end{center}

\begin{acquis}
\textbf{Commencez par T1.} C'est le seul test qui vous dira, en une exécution, l'étendue réelle du
travail --- et il est probable qu'il trouve des lectures d'environnement dont ni vous ni moi ne
soupçonnons l'existence. Mon propre inventaire est passé de sept à trente-deux en une commande.
\end{acquis}

\section{Annexe · Vérifier l'état de départ}

\renewcommand{\arraystretch}{1.22}
\begin{center}\small
\begin{tabular}{@{}>{\raggedright\arraybackslash}p{5.2cm} >{\raggedright\arraybackslash}p{10.8cm}@{}}
\toprule
\textbf{Affirmation} & \textbf{Ce qui la rejoue} \\
\midrule
25 variables dans la passerelle &
  \code{grep -rho 'EFE\_[A-Z0-9\_]*' crates/efe-swarm-gateway/src | sort -u | wc -l} \\
7 dans la console &
  \code{grep -rho 'EFE\_[A-Z0-9\_]*' crates/audit-console/src | sort -u | wc -l} \\
La discipline existe déjà &
  \code{sed -n '26,100p' crates/efe-swarm-gateway/src/main.rs} \\
Le garde est une fonction pure testée &
  \code{grep -n 'fn frontier\_boot\_guard' -A 6 crates/efe-swarm-gateway/src/main.rs} \\
\code{/readyz} fait une readiness réelle &
  \code{sed -n '1,45p' crates/efe-swarm-gateway/src/routes/health.rs} \\
Aucun document ne pose le profil &
  \code{grep -rl EFE\_LLM\_FRONTIER\_ENABLED . -{}-include='*.md' -{}-include='*.yml'} \sz{\fl{} vide} \\
\bottomrule
\end{tabular}
\end{center}

Un avertissement pour finir. Ce document décrit une \textbf{architecture}, pas un état : rien de ce
qu'il propose n'existe encore. Ce qui \emph{est} vérifiable aujourd'hui, ce sont les \textbf{trous} et
les \textbf{acquis} --- l'annexe donne la commande pour chacun. Le reste est une intention, et se
jugera au critère d'acceptation du §5.

\end{document}
